% Options for packages loaded elsewhere
\PassOptionsToPackage{unicode}{hyperref}
\PassOptionsToPackage{hyphens}{url}
\PassOptionsToPackage{dvipsnames,svgnames*,x11names*}{xcolor}
%
\documentclass[
  11pt,
]{article}
\usepackage{lmodern}
\usepackage{amssymb,amsmath}
\usepackage{ifxetex,ifluatex}
\ifnum 0\ifxetex 1\fi\ifluatex 1\fi=0 % if pdftex
  \usepackage[T1]{fontenc}
  \usepackage[utf8]{inputenc}
  \usepackage{textcomp} % provide euro and other symbols
\else % if luatex or xetex
  \usepackage{unicode-math}
  \defaultfontfeatures{Scale=MatchLowercase}
  \defaultfontfeatures[\rmfamily]{Ligatures=TeX,Scale=1}
\fi
% Use upquote if available, for straight quotes in verbatim environments
\IfFileExists{upquote.sty}{\usepackage{upquote}}{}
\IfFileExists{microtype.sty}{% use microtype if available
  \usepackage[]{microtype}
  \UseMicrotypeSet[protrusion]{basicmath} % disable protrusion for tt fonts
}{}
\makeatletter
\@ifundefined{KOMAClassName}{% if non-KOMA class
  \IfFileExists{parskip.sty}{%
    \usepackage{parskip}
  }{% else
    \setlength{\parindent}{0pt}
    \setlength{\parskip}{6pt plus 2pt minus 1pt}}
}{% if KOMA class
  \KOMAoptions{parskip=half}}
\makeatother
\usepackage{xcolor}
\IfFileExists{xurl.sty}{\usepackage{xurl}}{} % add URL line breaks if available
\IfFileExists{bookmark.sty}{\usepackage{bookmark}}{\usepackage{hyperref}}
\hypersetup{
  pdfauthor={KnackAU, Claude (Anthropic), Gemini (Google DeepMind)},
  colorlinks=true,
  linkcolor=Maroon,
  filecolor=Maroon,
  citecolor=Blue,
  urlcolor=Blue,
  pdfcreator={LaTeX via pandoc}}
\urlstyle{same} % disable monospaced font for URLs
\usepackage[margin=1in]{geometry}
\usepackage{color}
\usepackage{fancyvrb}
\newcommand{\VerbBar}{|}
\newcommand{\VERB}{\Verb[commandchars=\\\{\}]}
\DefineVerbatimEnvironment{Highlighting}{Verbatim}{commandchars=\\\{\}}
% Add ',fontsize=\small' for more characters per line
\newenvironment{Shaded}{}{}
\newcommand{\AlertTok}[1]{\textcolor[rgb]{1.00,0.00,0.00}{\textbf{#1}}}
\newcommand{\AnnotationTok}[1]{\textcolor[rgb]{0.38,0.63,0.69}{\textbf{\textit{#1}}}}
\newcommand{\AttributeTok}[1]{\textcolor[rgb]{0.49,0.56,0.16}{#1}}
\newcommand{\BaseNTok}[1]{\textcolor[rgb]{0.25,0.63,0.44}{#1}}
\newcommand{\BuiltInTok}[1]{#1}
\newcommand{\CharTok}[1]{\textcolor[rgb]{0.25,0.44,0.63}{#1}}
\newcommand{\CommentTok}[1]{\textcolor[rgb]{0.38,0.63,0.69}{\textit{#1}}}
\newcommand{\CommentVarTok}[1]{\textcolor[rgb]{0.38,0.63,0.69}{\textbf{\textit{#1}}}}
\newcommand{\ConstantTok}[1]{\textcolor[rgb]{0.53,0.00,0.00}{#1}}
\newcommand{\ControlFlowTok}[1]{\textcolor[rgb]{0.00,0.44,0.13}{\textbf{#1}}}
\newcommand{\DataTypeTok}[1]{\textcolor[rgb]{0.56,0.13,0.00}{#1}}
\newcommand{\DecValTok}[1]{\textcolor[rgb]{0.25,0.63,0.44}{#1}}
\newcommand{\DocumentationTok}[1]{\textcolor[rgb]{0.73,0.13,0.13}{\textit{#1}}}
\newcommand{\ErrorTok}[1]{\textcolor[rgb]{1.00,0.00,0.00}{\textbf{#1}}}
\newcommand{\ExtensionTok}[1]{#1}
\newcommand{\FloatTok}[1]{\textcolor[rgb]{0.25,0.63,0.44}{#1}}
\newcommand{\FunctionTok}[1]{\textcolor[rgb]{0.02,0.16,0.49}{#1}}
\newcommand{\ImportTok}[1]{#1}
\newcommand{\InformationTok}[1]{\textcolor[rgb]{0.38,0.63,0.69}{\textbf{\textit{#1}}}}
\newcommand{\KeywordTok}[1]{\textcolor[rgb]{0.00,0.44,0.13}{\textbf{#1}}}
\newcommand{\NormalTok}[1]{#1}
\newcommand{\OperatorTok}[1]{\textcolor[rgb]{0.40,0.40,0.40}{#1}}
\newcommand{\OtherTok}[1]{\textcolor[rgb]{0.00,0.44,0.13}{#1}}
\newcommand{\PreprocessorTok}[1]{\textcolor[rgb]{0.74,0.48,0.00}{#1}}
\newcommand{\RegionMarkerTok}[1]{#1}
\newcommand{\SpecialCharTok}[1]{\textcolor[rgb]{0.25,0.44,0.63}{#1}}
\newcommand{\SpecialStringTok}[1]{\textcolor[rgb]{0.73,0.40,0.53}{#1}}
\newcommand{\StringTok}[1]{\textcolor[rgb]{0.25,0.44,0.63}{#1}}
\newcommand{\VariableTok}[1]{\textcolor[rgb]{0.10,0.09,0.49}{#1}}
\newcommand{\VerbatimStringTok}[1]{\textcolor[rgb]{0.25,0.44,0.63}{#1}}
\newcommand{\WarningTok}[1]{\textcolor[rgb]{0.38,0.63,0.69}{\textbf{\textit{#1}}}}
\usepackage{longtable,booktabs}
% Correct order of tables after \paragraph or \subparagraph
\usepackage{etoolbox}
\makeatletter
\patchcmd\longtable{\par}{\if@noskipsec\mbox{}\fi\par}{}{}
\makeatother
% Allow footnotes in longtable head/foot
\IfFileExists{footnotehyper.sty}{\usepackage{footnotehyper}}{\usepackage{footnote}}
\makesavenoteenv{longtable}
\setlength{\emergencystretch}{3em} % prevent overfull lines
\providecommand{\tightlist}{%
  \setlength{\itemsep}{0pt}\setlength{\parskip}{0pt}}
\setcounter{secnumdepth}{-\maxdimen} % remove section numbering

\author{KnackAU, Claude (Anthropic), Gemini (Google DeepMind)}
\date{Shannon-Prime Project · 2026-05-19}

\begin{document}

\hypertarget{kste-the-friedman-sieve-system-and-implementation}{%
\section{KSTE \& The Friedman Sieve: System and
Implementation}\label{kste-the-friedman-sieve-system-and-implementation}}

\textbf{Paper IV --- Engineering the Friedman Stack on Snapdragon V69}

\emph{KnackAU, Claude (Anthropic), Gemini (Google DeepMind)}

\emph{Shannon-Prime Project · 2026-05-19}

\begin{center}\rule{0.5\linewidth}{0.5pt}\end{center}

\hypertarget{abstract}{%
\subsection{Abstract}\label{abstract}}

This paper specifies the system that realizes the Friedman Stack of
Paper III on existing Shannon-Prime infrastructure. The central
contribution is the \emph{Knight--Spinor Tree Encoder} (KSTE) --- the
exact, computable function
\(\Phi: \mathbb{R}^{128} \to \mathcal{T}_{60,3}\) that maps a continuous
Key vector to a 60-node, 3-label rooted tree, fitting inside the
existing 63-byte packed block. We give the full encoder kernel in
pseudocode and C, the homeomorphic-embedding decision kernel
(\(O(|T_Q| \cdot |T_K|)\) on CPU, fully vectorisable on Hexagon HVX via
predicate-register operations), and the Friedman-sieve integration into
the \texttt{sp\_hex\_sqfree\_cache\_t} write path. We then describe the
test plan, the parity gates, and the experimental schedule for the
perplexity / eviction-rate benchmarks that decide whether the framework
ships as default or remains an opt-in path.

\begin{center}\rule{0.5\linewidth}{0.5pt}\end{center}

\hypertarget{system-position}{%
\subsection{1. System Position}\label{system-position}}

Paper III's theory rests on KSTE. If the encoder cannot be implemented
cheaply, the theory is academic; if it can, the Friedman Sieve becomes a
drop-in replacement for variance-ranked top-K eviction with mathematical
guarantees attached.

This paper is the engineering specification. It assumes Papers I
(theory) and II (system) and Paper III (Friedman theory) as
prerequisites.

\hypertarget{block-format}{%
\subsection{2. Block Format}\label{block-format}}

The existing 63-byte Spinor block (Paper II §9.1) is reused without
modification:

\begin{verbatim}
Bytes  0– 27 : 14 × fp16 anchor coefficients (Knight skeleton)
Bytes 28– 58 : 60 lanes × (3-bit magnitude + 1-bit phase) = 31 bytes
Bytes 59– 62 : 4-byte amax (block-level scaling)
\end{verbatim}

The KSTE encoder constructs the tree from this exact byte layout. No
additional storage is required. The tree representation is
\emph{derived}, not stored: each block carries enough information to
reconstruct \(\Phi(K)\) on demand in \(O(60)\) operations.

For the \emph{evicted} trees (the ones we keep in cache as eviction
comparators), we add a separate 60-byte packed-tree representation:

\begin{Shaded}
\begin{Highlighting}[]
\KeywordTok{typedef} \KeywordTok{struct}\NormalTok{ \{}
    \DataTypeTok{uint8\_t}\NormalTok{ labels [}\DecValTok{15}\NormalTok{];   }\CommentTok{/* 60 nodes × 2 bits */}
    \DataTypeTok{uint8\_t}\NormalTok{ parents[}\DecValTok{45}\NormalTok{];   }\CommentTok{/* 60 nodes × 6 bits */}
\NormalTok{\} sp\_kste\_tree;            }\CommentTok{/* 60 bytes packed */}
\end{Highlighting}
\end{Shaded}

Total per-cache-slot footprint: 63 (existing) + 60 (tree) = \textbf{123
bytes} per K slot.

\hypertarget{the-kste-encoder}{%
\subsection{3. The KSTE Encoder}\label{the-kste-encoder}}

\hypertarget{high-level-pseudocode}{%
\subsubsection{3.1 High-level pseudocode}\label{high-level-pseudocode}}

\begin{verbatim}
function KSTE(block: sp_ok_q8_block_t) -> sp_kste_tree:
    # 1. Extract anchors and residuals from existing block
    anchors[0..13]   = fp16_decode(block.bytes[0..27])
    (mag[0..59], phase[0..59]) = unpack_residuals(block.bytes[28..58])
    amax             = fp32_decode(block.bytes[59..62])

    # 2. Compute order-invariant signature
    rank_a[k] = rank-of(|anchors[k]|) for k in 0..13
    sign_a[k] = sign(anchors[k])      for k in 0..13

    # 3. Build tree
    T = new_tree()
    root = T.add_root(label = ROOT)

    # 3a. Add anchor children, ordered by rank
    for k in argsort(rank_a, descending):
        if |anchors[k]| < tau_A * amax:
            break       # below threshold; stop
        T.add_child(root, label = A)

    # 3b. Attach residuals to dominating anchors
    for j in 0..59:
        if mag[j] == 0:
            continue
        # Phase determines label
        label = B if phase[j] == 0 else C
        # Find dominating anchor
        target_rank = rank-of-magnitude(mag[j], amax)
        k_star = first k such that rank_a[k] > target_rank * alpha
        parent = (T.children(root))[k_star] if k_star exists else root
        # Magnitude becomes depth: chain of length mag[j]
        current = parent
        for d in 0..(mag[j] - 1):
            current = T.add_child(current, label = label)

    # 4. Budget enforcement
    while T.node_count > 60:
        T.remove_weakest_leaf()  # leaf with lowest |residual|

    return T.pack()  # → sp_kste_tree (60 bytes)
\end{verbatim}

\hypertarget{c-reference-implementation}{%
\subsubsection{3.2 C reference
implementation}\label{c-reference-implementation}}

\begin{Shaded}
\begin{Highlighting}[]
\CommentTok{/* sp\_kste.h */}
\PreprocessorTok{\#define SP\_KSTE\_MAX\_NODES   60}
\PreprocessorTok{\#define SP\_KSTE\_TAU\_A       0.05f}
\PreprocessorTok{\#define SP\_KSTE\_ALPHA       0.7f}

\KeywordTok{typedef} \KeywordTok{enum}\NormalTok{ \{ SP\_KSTE\_ROOT = }\DecValTok{0}\NormalTok{, SP\_KSTE\_A = }\DecValTok{1}\NormalTok{, SP\_KSTE\_B = }\DecValTok{2}\NormalTok{, SP\_KSTE\_C = }\DecValTok{3}\NormalTok{ \} sp\_kste\_label;}

\KeywordTok{typedef} \KeywordTok{struct}\NormalTok{ \{}
    \DataTypeTok{uint8\_t}\NormalTok{ labels [}\DecValTok{15}\NormalTok{];   }\CommentTok{/* packed 2{-}bit labels */}
    \DataTypeTok{uint8\_t}\NormalTok{ parents[}\DecValTok{45}\NormalTok{];   }\CommentTok{/* packed 6{-}bit parent indices */}
    \DataTypeTok{uint8\_t}\NormalTok{ node\_count;}
    \DataTypeTok{uint8\_t}\NormalTok{ \_pad[}\DecValTok{3}\NormalTok{];}
\NormalTok{\} sp\_kste\_tree;            }\CommentTok{/* sizeof = 64 bytes (60 packed + 4 metadata) */}

\CommentTok{/* sp\_kste.c */}
\DataTypeTok{void}\NormalTok{ sp\_kste\_encode(}\DataTypeTok{const}\NormalTok{ sp\_ok\_q8\_block\_t *block, sp\_kste\_tree *out);}
\DataTypeTok{int}\NormalTok{  sp\_kste\_embed (}\DataTypeTok{const}\NormalTok{ sp\_kste\_tree *Q, }\DataTypeTok{const}\NormalTok{ sp\_kste\_tree *K);   }\CommentTok{/* 1 iff Q ⪯ K */}
\end{Highlighting}
\end{Shaded}

The encoder has no floating-point divisions; all magnitude comparisons
are integer ranks. The single fp16→fp32 step on the 14 anchors is the
only non-integer work and is bounded to \(O(14)\) per block.

\hypertarget{determinism-and-stability}{%
\subsubsection{3.3 Determinism and
stability}\label{determinism-and-stability}}

The encoder is \textbf{deterministic and order-invariant} by
construction. Two Key vectors with the same Knight-Skeleton-Spinor
signature produce the \emph{bit-identical} tree. The encoder is stable
under any rescaling of \(K\) that preserves the rank pattern and phase
pattern --- exactly the order-invariance of Paper III §4.1.

This means the encoder is also stable under Frobenius shimming (Paper I
§3.2): the \(\pi^k\) scale factor preserves ranks and signs, so
\(\Phi(\pi^k K) = \Phi(K)\). The KSTE layer is \emph{invariant} under
Frobenius --- a property that matters because the engine never has to
decide whether to encode before or after the shim.

\hypertarget{the-homeomorphic-embedding-kernel}{%
\subsection{4. The Homeomorphic-Embedding
Kernel}\label{the-homeomorphic-embedding-kernel}}

\hypertarget{decision-procedure}{%
\subsubsection{4.1 Decision procedure}\label{decision-procedure}}

The relation \(T_Q \preceq T_K\) is decidable in polynomial time. The
standard algorithm is the \emph{ordered tree embedding} dynamic program
(Kilpeläinen--Mannila, 1995), \(O(|T_Q| \cdot |T_K|^2)\) in the worst
case, with optimizations bringing typical-case complexity down to
\(O(|T_Q| \cdot |T_K|)\) when many partial embeddings fail early.

For 60-node trees the worst case is \(\approx 216{,}000\) operations,
easily inside one HVX dispatch.

\hypertarget{packed-bitwise-form}{%
\subsubsection{4.2 Packed bitwise form}\label{packed-bitwise-form}}

Both \(T_Q\) and \(T_K\) are stored as \texttt{sp\_kste\_tree} (60-byte
packed). The embedding test in vectorised form:

\begin{verbatim}
function sp_kste_embed_hvx(Q: tree, K: tree) -> {0,1}:
    # Compare label masks: for each Q label, find candidate K nodes
    # with matching label using vcmpeq predicate
    cand_A = vcmpeq(K.labels, A_mask)   # HVX 64-byte vector
    cand_B = vcmpeq(K.labels, B_mask)
    cand_C = vcmpeq(K.labels, C_mask)

    # For each Q node, mask candidate K nodes by label
    # then check parent-ancestor relation via parent-pointer transitive closure
    K_anc = transitive_closure(K.parents)   # precomputed once per K

    # Run the DP greedily under vector-predicate guard
    # ... see §4.3 for the full kernel
\end{verbatim}

The transitive-closure of parent pointers is computed once per K and
cached alongside the packed tree (an additional 60 × 60 / 8 = 450 bytes,
or rounded to a 512-byte block for alignment). The total cache slot
becomes 123 + 512 = \textbf{635 bytes per K slot} for the
embedding-accelerated form.

If memory is tight, the closure can be recomputed at test time at
\(O(60^2 / 8) = 450\) ops per K, well inside HVX budget.

\hypertarget{hexagon-kernel-sketch}{%
\subsubsection{4.3 Hexagon kernel sketch}\label{hexagon-kernel-sketch}}

\begin{Shaded}
\begin{Highlighting}[]
\CommentTok{/* sp\_hex\_kste\_embed.idl */}
\DataTypeTok{int}\NormalTok{ sp\_hex\_kste\_embed(}
    \DataTypeTok{const} \DataTypeTok{uint8\_t}\NormalTok{ *Q\_packed,   }\CommentTok{/* 60{-}byte tree */}
    \DataTypeTok{const} \DataTypeTok{uint8\_t}\NormalTok{ *K\_packed,   }\CommentTok{/* 60{-}byte tree */}
    \DataTypeTok{uint8\_t}\NormalTok{ *result            }\CommentTok{/* 0 or 1 */}
\NormalTok{);}

\CommentTok{/* sp\_hex\_kste\_embed\_imp.c — DSP{-}side */}
\DataTypeTok{int}\NormalTok{ sp\_hex\_kste\_embed(...) \{}
\NormalTok{    HVX\_Vector vQ\_labels   = *(HVX\_Vector *)Q\_packed;}
\NormalTok{    HVX\_Vector vK\_labels   = *(HVX\_Vector *)K\_packed;}

    \CommentTok{/* For each label class, build a vector predicate of K{-}candidates */}
\NormalTok{    HVX\_VectorPred pA = Q6\_Q\_vcmp\_eq\_VbVb(vK\_labels, vAconst);}
\NormalTok{    HVX\_VectorPred pB = Q6\_Q\_vcmp\_eq\_VbVb(vK\_labels, vBconst);}
\NormalTok{    HVX\_VectorPred pC = Q6\_Q\_vcmp\_eq\_VbVb(vK\_labels, vCconst);}

    \CommentTok{/* Walk Q in BFS order, maintaining the current "must descend from" frontier */}
    \CommentTok{/* Frontier represented as a 60{-}bit mask (one bit per K node) */}
    \DataTypeTok{uint64\_t}\NormalTok{ frontier = }\DecValTok{1}\BuiltInTok{ULL}\NormalTok{;   }\CommentTok{/* root */}
    \ControlFlowTok{for}\NormalTok{ (}\DataTypeTok{int}\NormalTok{ q = }\DecValTok{1}\NormalTok{; q \textless{} Q.node\_count; ++q) \{}
        \DataTypeTok{uint8\_t}\NormalTok{ q\_lbl   = unpack\_label(Q.labels, q);}
        \DataTypeTok{uint8\_t}\NormalTok{ q\_par   = unpack\_parent(Q.parents, q);}
        \DataTypeTok{uint64\_t}\NormalTok{ cand   = (q\_lbl == A) ? pA\_mask : (q\_lbl == B) ? pB\_mask : pC\_mask;}
        \CommentTok{/* Restrict cand to descendants of frontier */}
        \DataTypeTok{uint64\_t}\NormalTok{ reach  = expand\_descendants(frontier, K\_anc);}
        \DataTypeTok{uint64\_t}\NormalTok{ valid  = cand \& reach;}
        \ControlFlowTok{if}\NormalTok{ (valid == }\DecValTok{0}\NormalTok{) }\ControlFlowTok{return} \DecValTok{0}\NormalTok{;}
        \CommentTok{/* Pick lowest{-}rank candidate (smallest index in K\textquotesingle{}s BFS order) */}
        \DataTypeTok{int}\NormalTok{ chosen      = \_\_builtin\_ctzll(valid);}
\NormalTok{        frontier        = (}\DecValTok{1}\BuiltInTok{ULL}\NormalTok{ \textless{}\textless{} chosen);}
\NormalTok{    \}}
    \ControlFlowTok{return} \DecValTok{1}\NormalTok{;}
\NormalTok{\}}
\end{Highlighting}
\end{Shaded}

This is a greedy embedding; full correctness requires backtracking when
the greedy choice fails. The full algorithm uses a 60-element stack of
partial assignments and runs in \(\le 60 \times 60 = 3600\) steps
worst-case. The kernel fits in VTCM with room to spare.

\hypertarget{sieve-integration}{%
\subsection{5. Sieve Integration}\label{sieve-integration}}

\hypertarget{cache-structure}{%
\subsubsection{5.1 Cache structure}\label{cache-structure}}

\begin{Shaded}
\begin{Highlighting}[]
\KeywordTok{typedef} \KeywordTok{struct}\NormalTok{ \{}
\NormalTok{    sp\_ok\_q8\_block\_t k\_block;     }\CommentTok{/* 63 bytes — existing */}
\NormalTok{    sp\_kste\_tree     k\_tree;      }\CommentTok{/* 60 bytes — new */}
    \DataTypeTok{uint64\_t}\NormalTok{         k\_anc\_mask;  }\CommentTok{/* 8 bytes — transitive closure of parent pointers */}
    \CommentTok{/* total: 131 bytes per slot, 64{-}byte aligned to 192 bytes */}
\NormalTok{\} sp\_friedman\_slot\_t;}
\end{Highlighting}
\end{Shaded}

\hypertarget{write-path}{%
\subsubsection{5.2 Write path}\label{write-path}}

\begin{Shaded}
\begin{Highlighting}[]
\DataTypeTok{int}\NormalTok{ sp\_friedman\_cache\_write(}
\NormalTok{    sp\_friedman\_cache\_t *cache,}
    \DataTypeTok{const} \DataTypeTok{float}\NormalTok{ *K\_vec,        }\CommentTok{/* incoming Key */}
    \DataTypeTok{int}\NormalTok{ layer, }\DataTypeTok{int}\NormalTok{ head, }\DataTypeTok{int}\NormalTok{ pos}
\NormalTok{) \{}
    \CommentTok{/* 1. Existing pipeline: pack K into 63{-}byte Spinor block */}
\NormalTok{    sp\_ok\_q8\_block\_t new\_block;}
\NormalTok{    sp\_hex\_compress\_f32(K\_vec, head\_dim, \&new\_block);}

    \CommentTok{/* 2. KSTE: derive tree from block */}
\NormalTok{    sp\_kste\_tree new\_tree;}
\NormalTok{    sp\_kste\_encode(\&new\_block, \&new\_tree);}

    \CommentTok{/* 3. Sieve test: does new\_tree embed into ANY existing tree? */}
    \ControlFlowTok{for}\NormalTok{ (}\DataTypeTok{int}\NormalTok{ t = }\DecValTok{0}\NormalTok{; t \textless{} cache{-}\textgreater{}count; ++t) \{}
        \ControlFlowTok{if}\NormalTok{ (sp\_kste\_embed(\&new\_tree, \&cache{-}\textgreater{}slots[t].k\_tree)) \{}
            \CommentTok{/* Subsumed; evict (i.e., do nothing — token absorbed) */}
\NormalTok{            cache{-}\textgreater{}eviction\_count++;}
            \ControlFlowTok{return}\NormalTok{ SP\_FRIEDMAN\_EVICTED;}
\NormalTok{        \}}
\NormalTok{    \}}

    \CommentTok{/* 4. Novel; add to cache */}
    \ControlFlowTok{if}\NormalTok{ (cache{-}\textgreater{}count \textgreater{}= cache{-}\textgreater{}capacity) \{}
        \CommentTok{/* Cache full; evict by Knight{-}Skeleton variance (fallback) */}
\NormalTok{        sp\_friedman\_evict\_weakest(cache);}
\NormalTok{    \}}
    \DataTypeTok{int}\NormalTok{ slot = cache{-}\textgreater{}count++;}
\NormalTok{    cache{-}\textgreater{}slots[slot].k\_block = new\_block;}
\NormalTok{    cache{-}\textgreater{}slots[slot].k\_tree  = new\_tree;}
\NormalTok{    cache{-}\textgreater{}slots[slot].k\_anc\_mask = sp\_kste\_anc\_closure(\&new\_tree);}
    \ControlFlowTok{return}\NormalTok{ SP\_FRIEDMAN\_ADMITTED;}
\NormalTok{\}}
\end{Highlighting}
\end{Shaded}

\hypertarget{read-path}{%
\subsubsection{5.3 Read path}\label{read-path}}

Attention scoring proceeds as before --- the polynomial-ring kernel
reads \texttt{k\_block} and computes the dot product with the current
\texttt{Q}. The tree representation is \emph{write-only} from the
attention kernel's perspective: it is consulted only by the sieve.

This decoupling means the sieve is \emph{purely} an admission policy.
Attention scoring is unchanged. PPL effect is bounded by the
eviction-rate × per-evicted-token information content.

\hypertarget{experimental-schedule}{%
\subsection{6. Experimental Schedule}\label{experimental-schedule}}

The experiments fall into three tiers. Each tier has an explicit
ship/no-ship criterion.

\hypertarget{tier-1-encoder-sanity}{%
\subsubsection{6.1 Tier 1: Encoder sanity}\label{tier-1-encoder-sanity}}

\begin{longtable}[]{@{}lll@{}}
\toprule
\begin{minipage}[b]{0.18\columnwidth}\raggedright
Test\strut
\end{minipage} & \begin{minipage}[b]{0.24\columnwidth}\raggedright
Inputs\strut
\end{minipage} & \begin{minipage}[b]{0.49\columnwidth}\raggedright
Pass criterion\strut
\end{minipage}\tabularnewline
\midrule
\endhead
\begin{minipage}[t]{0.18\columnwidth}\raggedright
T1.1 --- Encoder is deterministic\strut
\end{minipage} & \begin{minipage}[t]{0.24\columnwidth}\raggedright
Same \(K\) × 1000 trials\strut
\end{minipage} & \begin{minipage}[t]{0.49\columnwidth}\raggedright
Bit-identical tree across all trials\strut
\end{minipage}\tabularnewline
\begin{minipage}[t]{0.18\columnwidth}\raggedright
T1.2 --- Encoder is order-invariant\strut
\end{minipage} & \begin{minipage}[t]{0.24\columnwidth}\raggedright
\(K\) vs \(\pi^k K\) × 100 trials\strut
\end{minipage} & \begin{minipage}[t]{0.49\columnwidth}\raggedright
Bit-identical tree\strut
\end{minipage}\tabularnewline
\begin{minipage}[t]{0.18\columnwidth}\raggedright
T1.3 --- Encoder is sign-respecting\strut
\end{minipage} & \begin{minipage}[t]{0.24\columnwidth}\raggedright
\(K\) vs \(-K\)\strut
\end{minipage} & \begin{minipage}[t]{0.49\columnwidth}\raggedright
Trees differ only in B↔C label swaps\strut
\end{minipage}\tabularnewline
\begin{minipage}[t]{0.18\columnwidth}\raggedright
T1.4 --- Budget enforced\strut
\end{minipage} & \begin{minipage}[t]{0.24\columnwidth}\raggedright
\(K \sim \mathcal{N}(0, I_{128})\) × 1000\strut
\end{minipage} & \begin{minipage}[t]{0.49\columnwidth}\raggedright
\(|V(\Phi(K))| \le 60\) always\strut
\end{minipage}\tabularnewline
\begin{minipage}[t]{0.18\columnwidth}\raggedright
T1.5 --- Anchor count matches\strut
\end{minipage} & \begin{minipage}[t]{0.24\columnwidth}\raggedright
Top-14 selection\strut
\end{minipage} & \begin{minipage}[t]{0.49\columnwidth}\raggedright
14 ± 2 children of root\strut
\end{minipage}\tabularnewline
\bottomrule
\end{longtable}

All five tests must pass before any sieve testing begins.

\hypertarget{tier-2-sieve-behaviour}{%
\subsubsection{6.2 Tier 2: Sieve
behaviour}\label{tier-2-sieve-behaviour}}

\begin{longtable}[]{@{}lll@{}}
\toprule
\begin{minipage}[b]{0.18\columnwidth}\raggedright
Test\strut
\end{minipage} & \begin{minipage}[b]{0.24\columnwidth}\raggedright
Inputs\strut
\end{minipage} & \begin{minipage}[b]{0.49\columnwidth}\raggedright
Pass criterion\strut
\end{minipage}\tabularnewline
\midrule
\endhead
\begin{minipage}[t]{0.18\columnwidth}\raggedright
T2.1 --- Sieve termination\strut
\end{minipage} & \begin{minipage}[t]{0.24\columnwidth}\raggedright
100k random tokens, encoder fixed\strut
\end{minipage} & \begin{minipage}[t]{0.49\columnwidth}\raggedright
Cache size \(\le \mathrm{antichain}(\mathcal{T}_{60,3})\)\strut
\end{minipage}\tabularnewline
\begin{minipage}[t]{0.18\columnwidth}\raggedright
T2.2 --- Eviction rate\strut
\end{minipage} & \begin{minipage}[t]{0.24\columnwidth}\raggedright
WikiText-103, ctx=2048\strut
\end{minipage} & \begin{minipage}[t]{0.49\columnwidth}\raggedright
\(\ge 20\%\) at steady state\strut
\end{minipage}\tabularnewline
\begin{minipage}[t]{0.18\columnwidth}\raggedright
T2.3 --- Sieve PPL ablation\strut
\end{minipage} & \begin{minipage}[t]{0.24\columnwidth}\raggedright
Gemma3-1B, sieve on/off\strut
\end{minipage} & \begin{minipage}[t]{0.49\columnwidth}\raggedright
\(\Delta\text{PPL} \le 0.5\%\)\strut
\end{minipage}\tabularnewline
\begin{minipage}[t]{0.18\columnwidth}\raggedright
T2.4 --- Refutation procedure\strut
\end{minipage} & \begin{minipage}[t]{0.24\columnwidth}\raggedright
Inject adversarial \(K\)\strut
\end{minipage} & \begin{minipage}[t]{0.49\columnwidth}\raggedright
Counterexample found in \(\le 10^4\) ops\strut
\end{minipage}\tabularnewline
\begin{minipage}[t]{0.18\columnwidth}\raggedright
T2.5 --- Closure-axiom invariant\strut
\end{minipage} & \begin{minipage}[t]{0.24\columnwidth}\raggedright
Cache snapshot, random subsets\strut
\end{minipage} & \begin{minipage}[t]{0.49\columnwidth}\raggedright
Big \(\cap\) Big = Big for all pairs\strut
\end{minipage}\tabularnewline
\bottomrule
\end{longtable}

T2.3 is the gating test: if PPL drift exceeds 0.5\%, the sieve is not
shippable as default. Tier 3 is contingent on T2.3 passing.

\hypertarget{tier-3-ultraproduct-attention}{%
\subsubsection{6.3 Tier 3: Ultraproduct
attention}\label{tier-3-ultraproduct-attention}}

\begin{longtable}[]{@{}lll@{}}
\toprule
\begin{minipage}[b]{0.18\columnwidth}\raggedright
Test\strut
\end{minipage} & \begin{minipage}[b]{0.24\columnwidth}\raggedright
Inputs\strut
\end{minipage} & \begin{minipage}[b]{0.49\columnwidth}\raggedright
Pass criterion\strut
\end{minipage}\tabularnewline
\midrule
\endhead
\begin{minipage}[t]{0.18\columnwidth}\raggedright
T3.1 --- Principal ⇒ Top-1\strut
\end{minipage} & \begin{minipage}[t]{0.24\columnwidth}\raggedright
Toy 16-token cache\strut
\end{minipage} & \begin{minipage}[t]{0.49\columnwidth}\raggedright
UltraAttn = Top-1 exactly\strut
\end{minipage}\tabularnewline
\begin{minipage}[t]{0.18\columnwidth}\raggedright
T3.2 --- Łoś on toy\strut
\end{minipage} & \begin{minipage}[t]{0.24\columnwidth}\raggedright
Hand-crafted properties\strut
\end{minipage} & \begin{minipage}[t]{0.49\columnwidth}\raggedright
Holds at limit iff \(U\)-large set\strut
\end{minipage}\tabularnewline
\begin{minipage}[t]{0.18\columnwidth}\raggedright
T3.3 --- PPL on LongBench\strut
\end{minipage} & \begin{minipage}[t]{0.24\columnwidth}\raggedright
UltraAttn vs softmax\strut
\end{minipage} & \begin{minipage}[t]{0.49\columnwidth}\raggedright
\(\Delta\text{PPL} \le 3\%\) either way\strut
\end{minipage}\tabularnewline
\begin{minipage}[t]{0.18\columnwidth}\raggedright
T3.4 --- RULER long-context\strut
\end{minipage} & \begin{minipage}[t]{0.24\columnwidth}\raggedright
UltraAttn at \(n \ge 32k\)\strut
\end{minipage} & \begin{minipage}[t]{0.49\columnwidth}\raggedright
Match or beat softmax\strut
\end{minipage}\tabularnewline
\begin{minipage}[t]{0.18\columnwidth}\raggedright
T3.5 --- Wall-time\strut
\end{minipage} & \begin{minipage}[t]{0.24\columnwidth}\raggedright
HVX kernel vs softmax\strut
\end{minipage} & \begin{minipage}[t]{0.49\columnwidth}\raggedright
Within 20\% wall-time\strut
\end{minipage}\tabularnewline
\bottomrule
\end{longtable}

Tier 3 is a research experiment. Even a \(+3\%\) PPL with \(-50\%\)
wall-time would be publishable; the framework's value is the new
attention primitive, not strict PPL.

\hypertarget{build-and-wire-up}{%
\subsection{7. Build and Wire-Up}\label{build-and-wire-up}}

\hypertarget{new-files}{%
\subsubsection{7.1 New files}\label{new-files}}

\begin{verbatim}
lib/shannon-prime/
  sp_kste.h          (encoder + embed API)
  sp_kste.c          (CPU reference impl)
  sp_kste_pack.c     (bit-packing helpers)

shannon-prime-engine/
  src/sp_friedman_cache.h
  src/sp_friedman_cache.cpp
  src/sp_ultraproduct_attn.h
  src/sp_ultraproduct_attn.cpp
  src/sp_hex_kste_embed.idl
  src/sp_hex_kste_embed_imp.c

tests/
  test_sp_kste.cpp
  test_sp_friedman_cache.cpp
  test_sp_ultraproduct_attn.cpp
\end{verbatim}

Approximate LOC budget: encoder 200, embedding 400, sieve integration
500, ultraproduct attention 600, tests 1500. \textbf{Total
\textasciitilde3200 LOC.}

\hypertarget{cmake-hooks}{%
\subsubsection{7.2 CMake hooks}\label{cmake-hooks}}

\begin{Shaded}
\begin{Highlighting}[]
\KeywordTok{option}\NormalTok{(SP\_FRIEDMAN\_SIEVE       }\StringTok{"Enable Friedman sieve eviction"}\NormalTok{ ON)}
\KeywordTok{option}\NormalTok{(SP\_KSTE\_ENCODER         }\StringTok{"Enable KSTE encoder for KV{-}trees"}\NormalTok{ ON)}
\KeywordTok{option}\NormalTok{(SP\_ULTRAPRODUCT\_ATTN    }\StringTok{"Enable ultraproduct attention path"}\NormalTok{ OFF)}
\end{Highlighting}
\end{Shaded}

\hypertarget{cli-flags}{%
\subsubsection{7.3 CLI flags}\label{cli-flags}}

\begin{Shaded}
\begin{Highlighting}[]
\ExtensionTok{sp{-}engine.exe}\NormalTok{ perplexity{-}sp }\KeywordTok{\textbackslash{}}
    \ExtensionTok{{-}{-}model}\NormalTok{ gemma3{-}1b.gguf }\KeywordTok{\textbackslash{}}
    \ExtensionTok{{-}{-}frobenius{-}quant}\NormalTok{ {-}p 41 {-}k 8 }\KeywordTok{\textbackslash{}}
    \ExtensionTok{{-}{-}poly{-}attn}\NormalTok{ {-}{-}ntt{-}crt }\KeywordTok{\textbackslash{}}
    \ExtensionTok{{-}{-}friedman{-}sieve} \KeywordTok{\textbackslash{}}
    \ExtensionTok{{-}{-}kste{-}tau{-}A}\NormalTok{ 0.05 {-}{-}kste{-}alpha 0.7 }\KeywordTok{\textbackslash{}}
    \ExtensionTok{{-}{-}ultraproduct{-}attn}\NormalTok{=principal }\KeywordTok{\textbackslash{}}
    \ExtensionTok{{-}{-}ctx}\NormalTok{ 4096 {-}{-}chunks 4}
\end{Highlighting}
\end{Shaded}

\hypertarget{environment-variables}{%
\subsubsection{7.4 Environment variables}\label{environment-variables}}

\begin{longtable}[]{@{}lll@{}}
\toprule
Variable & Default & Effect\tabularnewline
\midrule
\endhead
\texttt{SP\_FRIEDMAN\_ENABLE} & 0 & Enable sieve at
runtime\tabularnewline
\texttt{SP\_FRIEDMAN\_CAPACITY} & 4096 & Maximum cache slots after sieve
dedup\tabularnewline
\texttt{SP\_KSTE\_TAU\_A} & 0.05 & Anchor inclusion threshold (×
amax)\tabularnewline
\texttt{SP\_KSTE\_ALPHA} & 0.7 & Residual-to-anchor attachment
ratio\tabularnewline
\texttt{SP\_ULTRAPRODUCT\_MODE} & none & \texttt{none} /
\texttt{principal} / \texttt{nonprincipal}\tabularnewline
\bottomrule
\end{longtable}

\hypertarget{measurement-and-reporting}{%
\subsection{8. Measurement and
Reporting}\label{measurement-and-reporting}}

Each tier produces a structured report:

\begin{Shaded}
\begin{Highlighting}[]
\FunctionTok{\{}
  \DataTypeTok{"tier"}\FunctionTok{:} \DecValTok{2}\FunctionTok{,}
  \DataTypeTok{"test\_id"}\FunctionTok{:} \StringTok{"T2.3"}\FunctionTok{,}
  \DataTypeTok{"config"}\FunctionTok{:} \FunctionTok{\{}
    \DataTypeTok{"model"}\FunctionTok{:} \StringTok{"gemma3{-}1b"}\FunctionTok{,}
    \DataTypeTok{"ctx"}\FunctionTok{:} \DecValTok{2048}\FunctionTok{,}
    \DataTypeTok{"frobenius"}\FunctionTok{:} \StringTok{"p=41,k=8"}\FunctionTok{,}
    \DataTypeTok{"sieve"}\FunctionTok{:} \StringTok{"on"}
  \FunctionTok{\},}
  \DataTypeTok{"metrics"}\FunctionTok{:} \FunctionTok{\{}
    \DataTypeTok{"ppl\_with\_sieve"}\FunctionTok{:} \FloatTok{11.83}\FunctionTok{,}
    \DataTypeTok{"ppl\_without"}\FunctionTok{:}    \FloatTok{11.83}\FunctionTok{,}
    \DataTypeTok{"delta\_pct"}\FunctionTok{:}       \FloatTok{0.00}\FunctionTok{,}
    \DataTypeTok{"eviction\_rate"}\FunctionTok{:}   \FloatTok{0.34}\FunctionTok{,}
    \DataTypeTok{"tree\_size\_mean"}\FunctionTok{:}  \FloatTok{41.2}\FunctionTok{,}
    \DataTypeTok{"embed\_test\_us"}\FunctionTok{:}   \DecValTok{18}
  \FunctionTok{\},}
  \DataTypeTok{"verdict"}\FunctionTok{:} \StringTok{"PASS"}
\FunctionTok{\}}
\end{Highlighting}
\end{Shaded}

These JSON reports are committed alongside the engine; they form the
\emph{audit trail} for the framework's empirical claims.

\hypertarget{known-risks}{%
\subsection{9. Known Risks}\label{known-risks}}

\begin{enumerate}
\def\labelenumi{\arabic{enumi}.}
\item
  \textbf{Encoder discriminative resolution.} The KSTE encoder may not
  separate close-but-distinct semantic Keys. If T2.3 fails (PPL drift
  \textgreater{} 0.5\%), the encoder is the prime suspect; mitigations
  include increasing anchor count from 14 to 20, or extending labels to
  4-bit (a non-trivial block-format change).
\item
  \textbf{Eviction rate may be too high.} If the sieve evicts
  \textgreater{} 80\% of tokens, novel structure is being lost.
  Mitigation: tighten the embedding test (require \emph{strict}
  embedding rather than allowing trivial ones).
\item
  \textbf{Eviction rate may be too low.} If \textless{} 5\% of tokens
  are evicted, the sieve is doing no useful work. Mitigation: loosen the
  embedding test, or add a magnitude-similarity gate.
\item
  \textbf{HVX dispatch density.} The sieve adds one embedding test per
  write. At 32 layers × 8 heads × 4k tokens, that is 1M tests per
  inference pass. At 18 µs per HVX embed (T3.5 budget), that is 18
  seconds of compute. Mitigation: batched embedding tests (Strike 16
  analog).
\item
  \textbf{Determinism on shipped builds.} The KSTE encoder is
  deterministic \emph{given the calibration knight-mask}. Builds with
  different calibrations will produce different trees from the same
  Keys. Mitigation: ship the knight-mask with the model; treat it as a
  model artefact.
\end{enumerate}

\hypertarget{the-choice-operator-implementation}{%
\subsection{10. The Choice Operator
Implementation}\label{the-choice-operator-implementation}}

Paper III §11 introduces the axiomatic layer: a choice operator \(F\)
and the Extended-Domain Reduction axiom. This section specifies the
system-side implementation.

\hypertarget{the-canonical-total-order-prec_f}{%
\subsubsection{\texorpdfstring{10.1 The canonical total order
\(\prec_F\)}{10.1 The canonical total order \textbackslash prec\_F}}\label{the-canonical-total-order-prec_f}}

For determinism, the choice operator requires a fixed total order on
\(\mathcal{T}_{60,3}\). We use the \emph{packed lexicographic} order on
the 60-byte \texttt{sp\_kste\_tree} representation: compare the two byte
arrays component-wise, return the comparison of the first differing
byte. The order is trivially computable in \(O(60)\) operations,
deterministic across all platforms, and stable under the encoder's
invariants --- a tree and its bit-equivalent representation compare
equal.

\begin{Shaded}
\begin{Highlighting}[]
\CommentTok{/* sp\_kste.h */}
\DataTypeTok{int}\NormalTok{ sp\_kste\_compare(}\DataTypeTok{const}\NormalTok{ sp\_kste\_tree *a, }\DataTypeTok{const}\NormalTok{ sp\_kste\_tree *b);}
\CommentTok{/* returns \textless{}0 if a ≺\_F b, 0 if equal, \textgreater{}0 if a ≻\_F b */}
\end{Highlighting}
\end{Shaded}

\hypertarget{the-selector}{%
\subsubsection{10.2 The selector}\label{the-selector}}

\(F(A)\) is computed by linear scan with the comparator:

\begin{Shaded}
\begin{Highlighting}[]
\CommentTok{/* sp\_kste.h */}
\DataTypeTok{const}\NormalTok{ sp\_kste\_tree *sp\_kste\_select\_canonical(}
    \DataTypeTok{const}\NormalTok{ sp\_kste\_tree **candidates,}
    \DataTypeTok{int}\NormalTok{ n\_candidates}
\NormalTok{);}
\CommentTok{/* returns the ≺\_F{-}minimum tree in the candidate set */}
\end{Highlighting}
\end{Shaded}

Complexity: \(O(n \cdot 60)\) for \(n\) candidates. At cache size 4096,
the worst-case selection is \(\approx 245{,}000\) byte comparisons ---
well inside a single HVX dispatch and below the embedding-test budget.

\hypertarget{integration-with-the-sieve}{%
\subsubsection{10.3 Integration with the
sieve}\label{integration-with-the-sieve}}

The sieve's admission test (§5.2) currently returns a boolean. Under the
axiomatic layer, the admission decision is augmented:

\begin{Shaded}
\begin{Highlighting}[]
\KeywordTok{typedef} \KeywordTok{enum}\NormalTok{ \{}
\NormalTok{    SP\_FRIEDMAN\_EVICTED,}
\NormalTok{    SP\_FRIEDMAN\_REPLACED  }\CommentTok{/* admitted as F(A) for some matching class */}
\NormalTok{\} sp\_friedman\_decision;}
\end{Highlighting}
\end{Shaded}

The \texttt{REPLACED} path is the choice-operator's behaviour: when a
class of cached trees matches the new tree's structural query, the
engine selects \(F(A)\) and may replace the entries of \(A\) in the
cache with the canonical representative. This is a \emph{compaction}
operation --- the cache shrinks toward canonical witnesses over time,
bounded by the antichain count of \(\mathcal{T}_{60,3}\) from Theorem
3.2.

\hypertarget{the-extended-domain-reduction-unit-test}{%
\subsubsection{10.4 The Extended-Domain Reduction unit
test}\label{the-extended-domain-reduction-unit-test}}

The reduction axiom is a checkable invariant. For any structural
predicate \(\varphi\) and any selected \(v = F(A)\), the axiom asserts
\(\varphi(v) \Rightarrow \varphi^*(v)\), where \(\varphi^*\) is the
relativization of \(\varphi\) to the active-window subset of the cache.
We implement a small library of standard predicates and a top-level
check routine:

\begin{Shaded}
\begin{Highlighting}[]
\KeywordTok{typedef} \DataTypeTok{int}\NormalTok{ (*sp\_predicate\_t)(}\DataTypeTok{const}\NormalTok{ sp\_kste\_tree *T);}

\DataTypeTok{int}\NormalTok{ sp\_predicate\_anchor\_count(}\DataTypeTok{const}\NormalTok{ sp\_kste\_tree *T);}
\DataTypeTok{int}\NormalTok{ sp\_predicate\_label\_b\_count(}\DataTypeTok{const}\NormalTok{ sp\_kste\_tree *T);}
\DataTypeTok{int}\NormalTok{ sp\_predicate\_max\_depth(}\DataTypeTok{const}\NormalTok{ sp\_kste\_tree *T);}
\CommentTok{/* ... extensible per workload ... */}

\DataTypeTok{int}\NormalTok{ sp\_extended\_reduction\_check(}
    \DataTypeTok{const}\NormalTok{ sp\_kste\_tree *v,                }\CommentTok{/* canonical witness, v = F(A) */}
\NormalTok{    sp\_predicate\_t      phi,}
    \DataTypeTok{const}\NormalTok{ sp\_friedman\_cache\_t *cache\_RO   }\CommentTok{/* active window cache */}
\NormalTok{);}
\CommentTok{/* returns 1 iff phi(v) implies phi restricted to RO */}
\end{Highlighting}
\end{Shaded}

This routine is the engineering realization of the axiom. It runs a
finite scan of \texttt{cache\_RO}, evaluates \(\varphi\) on each entry,
and verifies the implication. It is primitive-recursive: no unbounded
search, no real-valued arithmetic. See \texttt{TEST-SUITE.md} T2.12 for
the gating test.

\hypertarget{performance-characteristics}{%
\subsubsection{10.5 Performance
characteristics}\label{performance-characteristics}}

The choice-operator layer adds a measurable, bounded cost per token:

\begin{itemize}
\tightlist
\item
  Canonical comparison: \(O(60)\) per pair.
\item
  Selection over candidate class of size \(k\): \(O(60 k)\).
\item
  Reduction check across \(|\mathrm{RO}|\):
  \(O(|\mathrm{RO}| \cdot \mathrm{cost}(\varphi))\).
\end{itemize}

At cache size 4096 and typical predicate cost \(O(60)\), the per-token
CPU overhead is \(\le 250\) µs --- well inside the existing wall-time
budget. On Hexagon HVX the comparison is parallelised to one vector
instruction per pair, dropping the cost by an order of magnitude. Total
axiomatic-layer overhead at production scale: \(\le 5\%\) of attention
wall-time.

\hypertarget{build-and-cli}{%
\subsubsection{10.6 Build and CLI}\label{build-and-cli}}

\begin{Shaded}
\begin{Highlighting}[]
\KeywordTok{option}\NormalTok{(SP\_CHOICE\_OPERATOR }\StringTok{"Enable axiomatic choice operator + ED Reduction"}\NormalTok{ ON)}
\end{Highlighting}
\end{Shaded}

\begin{Shaded}
\begin{Highlighting}[]
\ExtensionTok{sp{-}engine.exe}\NormalTok{ perplexity{-}sp }\KeywordTok{\textbackslash{}}
    \ExtensionTok{{-}{-}model}\NormalTok{ gemma3{-}1b.gguf }\KeywordTok{\textbackslash{}}
    \ExtensionTok{{-}{-}frobenius{-}quant}\NormalTok{ {-}p 41 {-}k 8 }\KeywordTok{\textbackslash{}}
    \ExtensionTok{{-}{-}poly{-}attn}\NormalTok{ {-}{-}ntt{-}crt }\KeywordTok{\textbackslash{}}
    \ExtensionTok{{-}{-}friedman{-}sieve} \KeywordTok{\textbackslash{}}
    \ExtensionTok{{-}{-}choice{-}operator} \KeywordTok{\textbackslash{}}
    \ExtensionTok{{-}{-}reduction{-}radius}\NormalTok{ 0   \# r in Paper III §11.4}\KeywordTok{;} \ExtensionTok{0}\NormalTok{ = strict consistency}
\end{Highlighting}
\end{Shaded}

\texttt{-\/-reduction-radius} is the fuzzy-class radius from Paper III
§11.4. The default \(r = 0\) enforces strict consistency; higher values
approach soft-attention behaviour without ever instantiating a softmax.
Per-workload defaults (code: \(r=0\); chat: \(r\) tuned) live in
\texttt{docs/CHOICE-OPERATOR-RADII.md}.

\hypertarget{conclusion}{%
\subsection{11. Conclusion}\label{conclusion}}

The Friedman Stack is implementable on existing Shannon-Prime
infrastructure with \textasciitilde3500 LOC of new code (including the
axiomatic-layer additions of §10) and zero new memory layouts on the
cache hot path. The encoder is order-invariant, deterministic, and
reuses the existing Knight-Skeleton/Spinor decomposition. The embedding
test fits in HVX with the same VTCM footprint as the existing
compress/decode pipeline. The sieve is purely an admission policy; the
choice operator adds a canonicality primitive that compacts the cache
toward \(\prec_F\)-minimal witnesses; the Extended-Domain Reduction
axiom is unit-testable as a primitive-recursive structural check.
Attention scoring remains the polynomial-ring + CRT-NTT kernel of Paper
II.

The experiments of §6 (extended by T2.12 and T3.6 in
\texttt{TEST-SUITE.md}) decide whether the framework ships. The
ship/no-ship gate is T2.3 --- a \(\le 0.5\%\) PPL drift with the sieve
enabled on Gemma3-1B at ctx=2048. If T2.3 passes, the Friedman Stack
becomes the default eviction policy in the next release; if not, it
remains an opt-in path while the encoder is iterated.

The next session continues with the implementation roadmap of
\texttt{IMPLEMENTATION-ROADMAP.md} and the test specifications of
\texttt{TEST-SUITE.md}.

\begin{center}\rule{0.5\linewidth}{0.5pt}\end{center}

\hypertarget{references}{%
\subsection{References}\label{references}}

\begin{enumerate}
\def\labelenumi{\arabic{enumi}.}
\tightlist
\item
  Kilpeläinen, P. \& Mannila, H., \emph{Ordered and Unordered Tree
  Inclusion}, SIAM J. Comput. 24, 1995.
\item
  Friedman, H. M., \emph{FOM Embedded Maximal Clique posts}, 2009--2018.
\item
  Łoś, J., \emph{Quelques remarques, théorèmes et problèmes sur les
  classes définissables d'algèbres}, 1955.
\item
  Qualcomm, \emph{Hexagon V69 HVX Intrinsics Reference}, 2023.
\item
  Shannon-Prime, \emph{Paper III --- The Friedman Stack} (theory, this
  work's companion).
\item
  Shannon-Prime, \emph{Paper I, II} (Frobenius framework + engine).
\item
  Anderson, C. A., \emph{Some Emendations on Gödel's Ontological Proof},
  1990.
\item
  Hilbert, D. \& Bernays, P., \emph{Grundlagen der Mathematik II}, 1939
  (epsilon-operator).
\item
  Ackermann, W., \emph{Zur Widerspruchsfreiheit der Zahlentheorie},
  Math. Ann. 117, 1940.
\end{enumerate}

\begin{center}\rule{0.5\linewidth}{0.5pt}\end{center}

\emph{This paper is consumed by \texttt{IMPLEMENTATION-ROADMAP.md}.
Tests are tracked in \texttt{TEST-SUITE.md}.}

\end{document}
